% Lösungen Einheit 1 von 3 – Schnittpunkt von Gerade und Ebene (zusammenbau v0.1)
\einheitenkopf{Einheit 1 von 3 · Schnittpunkt von Gerade und Ebene}

\erg{9}{a) ein Punkt – der Durchstoßpunkt.}
\erg{10}{a) einsetzen – die Geradenkoordinaten in die Koordinatengleichung. \quad b) Einsetzen: $(-1 + t) + 2 + (1 + t) = 6$, $2 + 2t = 6$, also $t = 2$; Schnittpunkt $S(1|2|3)$ \quad c) Richtungsvektor $\overrightarrow{PM} = (3|1|-3)$; $\vec{x} = (0|0|3) + t \cdot (3|1|-3)$ in $L$: $3t + 2t + 3 - 3t = 4$, also $t = 0{,}5$ und $Q(1{,}5|0{,}5|1{,}5)$ \quad d) $D$, $E$ und $F$ haben alle $z = 6$, die Ebene ist $z = 6$. Gerade: $\vec{x} = (-2|2|0) + t \cdot (-2|2|8)$; $8t = 6$ liefert $t = 0{,}75$ und den Punkt $(-3{,}5|3{,}5|6)$, also $R$. \quad e) Lichtrichtung $\overrightarrow{AA'} = (-2|1|-2)$; Lichtgerade $\vec{x} = (-2|4|4) + t \cdot (-2|1|-2)$; $4 - 2t = 0$ liefert $t = 2$ und $C'(-6|6|0)$ \quad f) Bahn $\vec{x} = (6|3|0) + r \cdot (-2|1|0)$, der Schritt $r = 1$ dauert 4 s; in $E$ eingesetzt ergibt sich $r = 1{,}5$; Zeit $r \cdot 4 = 6$ s \quad g) Die Sichtgrenze ist die Ebene durch das Auge und die obere Kante: Parameterform mit dem Auge als Stützpunkt und den Vektoren vom Auge zu den beiden Kantenendpunkten als Spannvektoren. Der gesuchte Punkt liegt zugleich auf der Kante: Geradengleichung durch ihre beiden Endpunkte. Beide Parameterformen gleichsetzen, das Gleichungssystem aus drei Gleichungen mit drei Parametern lösen und den Parameter der Kantengeraden in deren Gleichung einsetzen. Eine Gerade vom Auge über einen einzelnen Mauerpunkt reicht nicht.}
\erg{11}{a) Lichtgerade $\vec{x} = (0|0|h) + t \cdot (4|2|-4)$ trifft den Boden bei $t = \frac{h}{4}$: Schattenpunkt $(h|\frac{h}{2}|0)$; $(h - 1)^2 + \left(\frac{h}{2} + 2\right)^2 = 25$, also $1{,}25h^2 + 5 = 25$; die positive Lösung ist $h = 4$ m}
\erg{12}{a) Der Aufpunkt fehlt beim Einsetzen. Richtig: $(1 + t) + (2 + t) + 2t = 7$, also $3 + 4t = 7$, $t = 1$ und $S(2|3|2)$. \quad b) Jede Koordinate des Geradenpunkts hängt nur vom Parameter ab; in die eine Koordinatengleichung eingesetzt, bleibt eine Gleichung, in der nur der Parameter unbekannt ist.}
